% resumo em português
\setlength{\absparsep}{18pt} % ajusta o espaçamento dos parágrafos do resumo

\begin{resumo}
A disciplina de Teoria da Computação é fundamental na formação de discentes da área de tecnologia, fornecendo o embasamento teórico para o desenvolvimento de compiladores modernos. Contudo, conceitos relacionados às Linguagens Livres de Contexto (LLCs) e aos seus reconhecedores, os Autômatos de Pilha Não Determinísticos (APNs), frequentemente impõem barreiras ao aprendizado devido ao seu alto nível de abstração e à complexidade inerente de visualizar múltiplas computações paralelas. Este trabalho apresenta o desenvolvimento de uma ferramenta pedagógica baseada na web voltada à modelagem gráfica e à simulação em tempo real de APNs. A metodologia consistiu na implementação de uma aplicação executada integralmente no lado do cliente (\textit{frontend}), utilizando JavaScript puro e a biblioteca Cytoscape para a renderização de grafos. O sistema permite ao usuário desenhar autômatos, definir diferentes critérios de aceitação e analisar dinamicamente as ramificações geradas pelo não-determinismo por meio de estruturas denominadas \textit{snapshots}. Como resultado, a plataforma viabiliza a exportação e importação de modelos em formato JSON e o gerenciamento de exercícios de forma totalmente síncrona e gratuita. Embora o processamento em \textit{thread} única imponha limitações de desempenho frente a volumes massivos de dados, a ferramenta cumpre com sucesso o seu propósito didático de mitigar as dificuldades de visualização estática tradicionais. Como trabalhos futuros, propõem-se a realização de testes empíricos com os estudantes para validação da eficácia pedagógica, a expansão do escopo para outras máquinas de estados e a arquitetura de um \textit{backend} para otimização do desempenho.

\textbf{Palavras-chave}: Teoria da Computação. Autômatos de Pilha Não Determinísticos. Ferramenta Didática. Visualização de Algoritmos. Plataforma Web.
\end{resumo}